\clearpage
\VerdierPhysicalPageAnchor{0208}%
\begingroup\small\setlength{\parskip}{0pt}\setlength{\emergencystretch}{1.5em}
\setlength{\abovedisplayskip}{4pt plus 2pt minus 2pt}%
\setlength{\belowdisplayskip}{4pt plus 2pt minus 2pt}%

\VerdierSemanticAnchor{ch3d-3}
\section{Étude des \texorpdfstring{$\operatorname{Ext}$}{Ext}.}

\VerdierSemanticAnchor{ch3d-3-preamble}
\noindent\textit{Dans ce paragraphe, $\mathcal A$ désigne une catégorie
abélienne.}

\VerdierSemanticAnchor{ch3d-3-1}
\subsection{Différentes définitions équivalentes.}

\VerdierSemanticAnchor{ch3d-3-1-1}
\noindent\textbf{3.1.1.}\quad
On dit que la \emph{dimension cohomologique de $\mathcal A$ est inférieure
ou égale à $n$} si pour tout couple $X$ et $Y$ d'objets de $\mathcal A$ et
pour tout entier $p>n$, on a
$\operatorname{Ext}_{\mathcal A}^{p}(X,Y)=0$ ((1.2.5), (1.2.9)). On appelle
\emph{dimension cohomologique de $\mathcal A$} et on note
$\operatorname{dimcoh}(\mathcal A)$, le plus petit des entiers $n$ pour
lesquels la dimension cohomologique de $\mathcal A$ est inférieure ou égale
à $n$. S'il n'existe pas d'entier $n$ pour lequel la dimension cohomologique
de $\mathcal A$ soit inférieure ou égale à $n$, on dit que la dimension
cohomologique de $\mathcal A$ est infinie et on écrit
$\operatorname{dimcoh}(\mathcal A)=+\infty$.

\VerdierSemanticAnchor{ch3d-3-1-2}
\noindent\textbf{3.1.2.}\quad
Rappelons (1.1.3) qu'on désigne par $\operatorname{Qis}^{*}$ l'ensemble des
quasi-isomorphismes de $K^{*}(\mathcal A)$
($*=b,-,+,{}$ ``vide''). Soit $X$ un objet de $K^{*}(\mathcal A)$, on
désigne par $\operatorname{Qis}^{*}/X$ (resp.\
$X\backslash\operatorname{Qis}^{*}$) la catégorie des quasi-isomorphismes de
but $X$ (resp.\ de source $X$) (chap. II, 2.2.1). L'ensemble
$\operatorname{Qis}^{*}$ étant un système multiplicatif de morphismes, la
catégorie $\operatorname{Qis}^{*}/X$ (resp.\
$X\backslash\operatorname{Qis}^{*}$) est filtrante à gauche (resp.\ à droite)
(chap. II, 2.2.1). La catégorie $\operatorname{Qis}^{*}/X$ (resp.\
$X\backslash\operatorname{Qis}^{*}$) admet un objet initial (resp.\ final)
$u\colon Y\longrightarrow X$ (resp.\ $u\colon X\longrightarrow Y$) si et
seulement s'il existe un objet $Q^{*}$-libre à gauche (resp.\ à droite) $Y$ et
un quasi-isomorphisme $u\colon Y\longrightarrow X$ (resp.\
$u\colon X\longrightarrow Y$), comme on le voit immédiatement en utilisant
(chap. II, 2.3.3).

\VerdierSemanticAnchor{ch3d-3-1-3}
\medskip\noindent\textbf{Proposition 3.1.3.}\quad
\textit{Soient $X$ et $Y$ deux complexes de $\mathcal A$.}

\VerdierSemanticAnchor{ch3d-3-1-3-1}
\[
\textit{(1)}\qquad
\operatorname{Ext}_{\mathcal A}^{0}(X,Y)=
\varinjlim_{(\operatorname{Qis}/X)^o}
\operatorname{Hom}_{K(\mathcal A)}(\mathord\cdot,Y).
\]

\VerdierSemanticAnchor{ch3d-3-1-3-2}
\[
\textit{(2)}\qquad
\operatorname{Ext}_{\mathcal A}^{0}(X,Y)\simeq
\varinjlim_{Y\backslash\operatorname{Qis}}
\operatorname{Hom}_{K(\mathcal A)}(X,\mathord\cdot).
\]

\VerdierSemanticAnchor{ch3d-3-1-3-3}
\[
\textit{(3)}\qquad
\operatorname{Ext}_{\mathcal A}^{0}(X,Y)\simeq
\varinjlim_{(\operatorname{Qis}/X)^o\times
  (Y\backslash\operatorname{Qis})}
\operatorname{Hom}_{K(\mathcal A)}(\mathord\cdot,\mathord\cdot).
\]

\VerdierSemanticAnchor{ch3d-3-1-3-4}
\[
\textit{(4)}\qquad
\operatorname{Ext}_{\mathcal A}^{n}(X,Y)\simeq
\operatorname{Ext}_{\mathcal A}^{0}(X[p],Y[n+p]),
\qquad\textit{pour tout entier $p$.}
\]

\VerdierSemanticAnchor{ch3d-3-1-3-5}
\[
\textit{(5) Pour tout $X$ dans $K^{-}(\mathcal A)$,}\qquad
\operatorname{Ext}_{\mathcal A}^{0}(X,Y)\simeq
\varinjlim_{(\operatorname{Qis}^{-}/X)^o}
\operatorname{Hom}_{K(\mathcal A)}(\mathord\cdot,Y).
\]

\VerdierSemanticAnchor{ch3d-3-1-3-6}
\[
\textit{(6) Pour tout $Y$ dans $K^{+}(\mathcal A)$,}\qquad
\operatorname{Ext}_{\mathcal A}^{0}(X,Y)\simeq
\varinjlim_{Y\backslash\operatorname{Qis}^{+}}
\operatorname{Hom}_{K(\mathcal A)}(X,\mathord\cdot).
\]
\endgroup
